% SOURCE_STATEMENT: "For the Edwards-Anderson Ising spin glass on $\\mathbb{Z}^d$ with i.i.d. mean-zero finite-variance nearest-neighbor couplings, prove or disprove the existence of a thermodynamic phase transition above a lower critical dimension $d_c$, determine $d_c$, decide whether a low-temperature phase breaks spin-flip symmetry, and characterize the number and metastate organization of pure-state pairs."
% Preparation and final self-audit: CONTRIBUTING.md prompts C and D; source/audit artifacts retained in the local physics_statements_2026_10_11 batch.
% !TeX program = lualatex
\documentclass[11pt]{article}
\usepackage[margin=25mm]{geometry}
\usepackage{fontspec}
\setmainfont{Latin Modern Roman}
\usepackage{amsmath,amssymb,mathtools}
\usepackage{unicode-math}
\setmathfont{Latin Modern Math}
\usepackage{xurl,hyperref}
\hypersetup{colorlinks=true,urlcolor=blue}
\setcounter{secnumdepth}{0}
\setlength{\parindent}{0pt}
\setlength{\parskip}{5pt}
\setlength{\emergencystretch}{3em}
\newcommand{\sref}[2]{\hyperlink{source:#1:#2}{[S#2]}}
\newcommand{\cataloguescope}[1]{\paragraph{Recorded target.}#1}
\newcommand{\R}{\mathbb R}
\newcommand{\C}{\mathbb C}
\newcommand{\N}{\mathbb N}
\newcommand{\Z}{\mathbb Z}
\newcommand{\Q}{\mathbb Q}
\newcommand{\F}{\mathbb F}
\newcommand{\eps}{\varepsilon}
\begin{document}
\section{Short-Range Spin Glasses}\label{statement}

\subsection{Definitions and mathematical statement}

\paragraph{Shared notation.}$\mathbb N=\{1,2,\ldots\}$; the standard number fields have their usual meanings. All additional conventions are stated locally.


\paragraph{Definitions and notation.}In the Edwards--Anderson model, $\sigma_x\in\{-1,1\}$ and $H_\Lambda(\sigma)=-\sum_{\langle x,y\rangle}J_{xy}\sigma_x\sigma_y$ at zero external field. The nearest-neighbour $J_{xy}$ are independent and identically distributed, with mean zero and finite variance. Infinite-volume Gibbs measures satisfy the conditional Gibbs laws; pure states are their extreme points. Spin flip sends all spins to their negatives. A metastate is a probability measure on Gibbs measures, encoding possible volume-dependent limiting states.

\paragraph{Catalogue wording retained for traceability.}
For almost every coupling realization on $\mathbb Z^d$, prove or disprove a thermodynamic transition above a lower critical dimension $d_c$, and determine $d_c$ if a transition exists. Decide whether some low-temperature phase breaks spin-flip symmetry. For each applicable $d,T$, determine the number of pure-state pairs. If there are infinitely many, establish or refute the source’s proposed metastate picture: infinitely many mixed states, each the equal-weight mixture of a single spin-flip pair, with the visible pair changing with volume. If multiple Gibbs states exist, determine whether that phase persists under a sufficiently small external magnetic field.
\paragraph{Scope and source qualification.}The final small-field question and the almost-sure quantifier are restored from the original source; they were omitted from the short catalogue summary.
\subsection{Short English statement}
Determine whether short-range spin glasses have a phase transition, how many pure-state pairs occur, and how those pairs appear in large volumes.
\subsection{Sources}
\begin{itemize}
\item[\hypertarget{source:statement:1}{S1}] ULAM.AI and the UnsolvedMath Contributors, UnsolvedMath, catalogue entry. Catalogue curation under CC BY 4.0; definitions and English exposition added locally.
\url{https://huggingface.co/datasets/ulamai/UnsolvedMath}.
\item[\hypertarget{source:statement:2}{S2}] C. M. Newman and D. L. Stein, original contribution to Open Problems in Mathematical Physics (1998), edited by Michael Aizenman.
\url{https://web.math.princeton.edu/~aizenman/OpenProblems_MathPhys/9803.SpinGlass.html}.
\end{itemize}
\label{end:statement}
\end{document}
